\documentclass[11pt]{article}
\usepackage[T1]{fontenc}
\usepackage{lmodern}
\usepackage[margin=1in]{geometry}
\usepackage{graphicx}
\usepackage{amsmath}
\usepackage{microtype}
\usepackage{placeins}
\usepackage[font=small,skip=6pt]{caption}
\usepackage[numbers,sort&compress]{natbib}
\usepackage[hidelinks]{hyperref}

\graphicspath{{results/chip_metabolic_graph/}}
\setcounter{topnumber}{4}
\setcounter{bottomnumber}{4}
\setcounter{totalnumber}{4}
\renewcommand{\topfraction}{0.9}
\renewcommand{\bottomfraction}{0.9}
\renewcommand{\textfraction}{0.05}
\renewcommand{\floatpagefraction}{0.7}

\title{Metabolomics of Tet2 Knockout in Clonal Haematopoiesis\\
of Indeterminate Potential (CHIP)}
\author{}
\date{}

\begin{document}
\maketitle

\begin{abstract}
    Age-associated myelopoiesis is the myeloid demand of aged marrow. Clonal
    Haematopoiesis of Indeterminate Potential (CHIP) clones accumulate with years and
    turn interleukin-1 (IL-1)-driven emergency myelopoiesis into a chronic state.
    Previously pusblished work mapped TET2 and DNMT3TA loss to myeloid lipid anabolism and OGT, O-GlcNAc and ACLY lipid circuit in myeloid cells. Polar
    metabolites from HSPCs under IL-1 have not been measured. We now have designed a model of to simulate this 2$\times$2 experiment. We used a visible neural network (VNN) to model the metabolic changes that occur when one gene-metabolomic node differences between groups. Gene marker plots highlight where the encding intreaction is possible. This done by testing whether the mean gene-embedding interaction in its gene set is more extreme than random gene sets of the same size over 9{,}999 shuffles. The $p$ under each 2$\times$2 panel is that value. Mouse HSPCs meet it for HMP shunt ($p=0.0001$) and human HSPCs meet it for glycolysis, Complex~I and Complex~II (all $p=0.0001$), HMP shunt ($p=0.0014$), Krebs-cycle NADH generation
    ($p=0.006$), and ATP regeneration from glucose ($p=0.043$). This suggests a proposed experiment of measuring the polar metabolites from the HSPCs under IL-1 conditions and validate the model either in cells or medium. This will be done by first performing a cheap screens to test whether the cargo is used for output, catalysis, or PPP fitness. Then performing polar liquid chromatography mass spectrometry (LC-MS) of cells
    and medium to test whether lactate, succinate, itaconate, and $\alpha$-KG leave the cell.
\end{abstract}

\section*{Introduction}

CHIP is an expanded blood clone carrying a myeloid-malignancy driver at variant allele
frequency of at least 2\%, without unexplained cytopenia or frank neoplasm. That
threshold is a sequencing convention. TET2 and DNMT3A are
the common drivers. Clones accumulate with years. Age-associated myelopoiesis is the
myeloid demand of aged marrow, and interleukin-1 (IL-1) is a main cytokine of that tone.
IL-1 on young wild-type haematopoietic stem and progenitor cells (HSPCs) is emergency
myelopoiesis as an acute event. CHIP turns that event-programme chronic. Mutant myeloid
cells then amplify inflammasome tone, the clone expands, wild-type haematopoietic
stem cells (HSCs) exhaust, and inflammaging is amplified. We start with Tet2 knockout
because TET2 is the metabolic clone. DNMT3A is the clone with stronger bone and frailty
evidence; it is not this polar panel. Kim, Hergott and Ebert mapped TET2-mutant innate
activation onto O-GlcNAc and ACLY-linked lipid anabolism in myeloid
cells~\cite{Kim2026TET2metabolism}. Gene-metabolite task scores already
indicate which polar nodes differ between Tet2 and wild-type under IL-1.
Those indicated metabolites are the hypothesis this experiment measures in the
osteoimmune niche that inflammation occupies.

In marrow, IL-1 expands Tet2-mutant HSPCs and biases them toward myeloid output. Human TET2-edited HSPCs under
LPS ask the same question in another species. Gene-level RNA on that 2$\times$2 indicates
which metabolic nodes differ from a shared gene profile.
Metabolic-task scores sit downstream of transcript counts and are not flux. The
hypothesis is that Tet2 HSPCs export lactate, succinate, itaconate, and $\alpha$-KG under
IL-1 at levels that differ from wild-type on the same 2$\times$2. Cheap cell assays on
the edited HSPCs come first, to de-risk that measurement. Colony assays,
5-hydroxymethylcytosine, and G6PD or Seahorse test whether that RNA indication is used
for output, catalysis, or PPP fitness. Polar LC-MS of HSPC cells and medium is still
planned on the same cultures. It asks whether the indicated metabolites leave the cell.
A later SUCNR1 assay would test whether medium succinate is sensed. Polar metabolomics
extends the RNA indication into metabolites the niche, bone, and clone can sense.

\section*{Background}

Bone marrow is a shared osteohaematopoietic niche. Haematopoietic stem cells, stroma,
vessels, osteoblasts, osteoclasts, macrophages and T cells exchange cytokines, vesicles
and metabolites~\cite{He2026ImmuneAgeingSkeletal,Winter2024CHOsteoNiche}.
Inflammaging is the low-grade, chronic inflammatory tone that accumulates with
years~\cite{Franceschi2014Inflammaging}. Frailty is vulnerability to stressors as
chronic inflammation weakens reserve.

Emergency myelopoiesis is the demand-adapted shift of HSPCs toward rapid myeloid output
during an event~\cite{Park2024IL1myelopoiesis,Pietras2016IL1HSC}. IL-1 on young
wild-type HSPCs is that event. It is an acute pulse, not aging.

Age-associated myelopoiesis is the myeloid demand of aged marrow. Older marrow is
myeloid-biased. CHIP-negative marrow can still be myeloid-biased. A marrow CD4 CTL
$\to$ CCL5--CCR5 instructor is one such path~\cite{GabandeRodriguez2026CD4CTL}. CHIP
clones accumulate with years. In people, CHIP turns the emergency programme chronic.
Mature mutant myeloid cells over-secrete the cytokines that gave them a competitive
edge, notably IL-1$\beta$ from TET2-deficient
macrophages~\cite{Fuster2017TET2athero,Jaiswal2017CVD,Fuster2020TET2insulin}. Marrow
remains in an emergency-myelopoiesis state, and the mutated clone expands
further~\cite{Caiado2022Tet2IL1,Caiado2024IL1AgingHSC,McClatchy2023TET2IL1,
    Hajishengallis2026CHIPinflammaging}.
Wild-type HSCs under that continuous demand suffer proliferative exhaustion and impaired
regeneration~\cite{Pietras2016IL1HSC,Kovtonyuk2022IL1Microbiome}. Once the clone is
large, its myeloid daughters amplify systemic inflammaging with the same
cytokines~\cite{Franceschi2014Inflammaging,Caiado2024IL1AgingHSC,
    Motisi2026CHIPInflammageing}.
Fuster and colleagues already made that myeloid IL-1$\beta$ cargo causal for insulin
resistance in ageing and obesity~\cite{Fuster2020TET2insulin}. That is the CHIP limb.
It is not chronological aging.

Sorted HSPCs do not close the macrophage IL-1 loop, so a young Tet2 $\times$ IL-1
2$\times$2 still adds the cytokine. The four arms are a clone you can knock out plus
IL-1 you add. Polar LC-MS is the readout. Age-associated myelopoiesis is why that loop
exists in people. It is not what the four arms measure. No public metabolome has
measured those species from Tet2 HSPCs under IL-1. Polar export of that panel into the
niche that already carries IL-1 is the cargo that has not been measured.

He and colleagues synthesise skeletal fragility as a composite of that aged osteoimmune
niche, not bone mineral density (BMD) alone~\cite{He2026ImmuneAgeingSkeletal}. CHIP-
associated myeloid inflammatory amplification is one limb, shaped by driver and clone
size. Direct bone evidence is strongest for DNMT3A-mutant
clones~\cite{Kim2021Dnmt3aOsteoporosis,Wang2024DNMT3Abone}. Bone-specific evidence for
TET2 remains limited~\cite{He2026ImmuneAgeingSkeletal,Winter2024CHOsteoNiche}.

DNMT3A models already explain inflammatory bone loss. Polar cargo is the TET2
metabolic extension, not a DNMT3A bone assay. Gene-metabolite task scores
already indicate between-group differences at succinate, $\alpha$-KG, lactate, and PPP
nodes. No polar metabolome has measured those metabolites from Tet2 HSPCs under IL-1.
That cargo reaches the osteoimmune niche He places under skeletal fragility. It also
reaches IL-1-driven clone expansion that feeds inflammaging. Extracellular succinate binds SUCNR1 on osteoclast-lineage cells and
drives osteoclastogenesis, a bone-remodeling path into frailty~\cite{Guo2017Succinate}.
The same receptor on HSPCs restricts haematopoiesis. Succinate can expand clones when
SUCNR1 is low, an autocrine path into the CHIP-inflammaging
loop~\cite{Cuminetti2026SUCNR1}. 4-Octyl itaconate alkylates GAPDH and is
anti-inflammatory in vivo, adjacent to resolution in that same tone~\cite{Liao20194OI}.
TET2 uses $\alpha$-KG as co-substrate. Succinate, fumarate and 2-hydroxyglutarate
compete with it~\cite{Zhang2025TET2dioxygenase}. Ascorbate is high in HSCs and feeds
TET2~\cite{Agathocleous2017Ascorbate}. A polar LC-MS panel on this HSPC 2$\times$2 would
test whether those cofactors and competitors move with Tet2 under IL-1.

The TET2 metabolic paper this extension is written against is Kim, Hergott, Ebert and
colleagues~\cite{Kim2026TET2metabolism}. An in vivo genome-wide CRISPR screen in primary
wild-type and Tet2-knockout HSPCs, read out in zymosan peritonitis, recovered Ogt and
Acly as Tet2-selective inflammation nodes. Tet2 restrains O-linked $N$-acetylglucosamine
transferase (OGT). Loss raises O-GlcNAc, H3K4me3 at lipid-anabolic loci, and ATP-citrate
lyase (ACLY) activity. Mouse neutrophils, bone-marrow-derived macrophages (BMDMs), and
human CD34$^{+}$-derived macrophages accumulate di- and triacylglycerols and cholesterol
esters. Polar metabolite profiling of Tet2-knockout BMDMs ranked glycerophospholipid
metabolism, not lactate, succinate, itaconate, or $\alpha$-KG export. ACLY inhibition
blunts Tet2-knockout hyperinflammation, including Il1b. That is myeloid lipid anabolism
under innate challenge. It is not FACS-sorted HSPC polar LC-MS on a Tet2 $\times$ IL-1
2$\times$2, and it is not conditioned-medium cargo into the osteoimmune niche.

Ni{\~n}o and colleagues report that Tet2 HSPCs overexpress glycolysis and OXPHOS genes,
enlarge mitochondria, and need glucose-6-phosphate dehydrogenase (G6PD) and the PPP to
keep redox and clone fitness~\cite{Nino2026TETmetabolism}. Human TET2 splits by lineage.
HSPCs blunt AP-1 inflammatory adaptation. Monocyte-macrophage daughters amplify
it~\cite{HuergaEncabo2026TET2inflammation}. Polar cargo from HSPCs can occupy the
inflammaging niche without mixing those daughters into the assay.

No public metabolome has measured those species from Tet2 HSPCs under IL-1.
McClatchy GSE209994 and Huerga Encabo GSE285379 are
RNA-seq~\cite{McClatchy2023TET2IL1,HuergaEncabo2023TET2neutrophil}. Agathocleous and
Lalioti supply the low-input polar LC-MS
method~\cite{Agathocleous2017Ascorbate,Lalioti2025HSCmetabolome}. Kim supplies the TET2
myeloid lipidome, and Ni{\~n}o supplies the HSPC fitness map. Measuring lactate,
succinate, itaconate, and $\alpha$-KG/succinate from Tet2 HSPCs under IL-1 places those
priors, and the gene-metabolite indications, in the inflammaging niche.

\section*{Approach}

We trained a metabolic visible neural network (VNN) on scCellFie gene scores for
glycolysis, oxidative phosphorylation/TCA, and pentose-phosphate (PPP) tasks in young
McClatchy GSE209994 marker-HSPCs (WT/Tet2 $\times$ vehicle/IL-1, 8
libraries)~\cite{McClatchy2023TET2IL1,Armingol2025scCellFie}. A VNN is a neural net
whose layers follow a knowledge graph rather than a dense hidden stack. Mouse and human
each get their own VNN. Each is trained for 100 epochs on GPU, reconstructing every cell
in that species' 2$\times$2. The network reconstructs each cell's gene scores through
that hierarchy, then the task encodings are frozen. The final figure is the factorial
interaction map, one 2$\times$2 per metabolic task from the frozen encodings
(Fig.~\ref{fig:mice}). The statistic is
$(\mathrm{Tet2_{IL1}}-\mathrm{WT_{IL1}})-(\mathrm{Tet2_{vehicle}}-\mathrm{WT_{vehicle}})$.
Colour is the four-arm mean, centered and scaled so the largest absolute arm is 1
representing the relative mean of the interaction. Panel titles are metabolic tasks and
match the gene-cloud labels. Those panels are VNN encodings of treatment within
genotype. Relative mean can look strong when the encodings are small. Glycolysis has
the most positive Tet2/IL-1 cell in that scaled map. Complex~I has the most negative.

Task-gene embeddings on that same 2$\times$2 sit in one pile: the genes share one
shift, so the encoding interaction between groups is not a usable test. Each task is
tested by a competitive mean permutation: whether the mean frozen gene-embedding
interaction in that gene set is more extreme than random sets of the same size. The $p$
under each panel is $(1+k)/10000$ from 9{,}999 gene shuffles, where $k$ is how many
shuffles were at least as extreme. Bonferroni over seven mouse tasks or eight human
tasks is about $0.006$ to $0.007$; 9{,}999 shuffles put the floor at $0.0001$, below
that cutoff. Bold titles mark $p<0.05$.

Human GSE285379 FACS-HSPCs (TET2 $\times$ LPS, 4 libraries) use a separate VNN of the
same hierarchy (Fig.~\ref{fig:human})~\cite{HuergaEncabo2023TET2neutrophil}. The
competitive mean test meets $p<0.05$ for glycolysis, Complex~I and Complex~II (all
$p=0.0001$), HMP shunt ($p=0.0014$), Krebs-cycle NADH generation ($p=0.006$), and ATP
regeneration from glucose ($p=0.043$). Bonferroni $0.05/8$ keeps the first five except
ATP regeneration. Pyruvate decarboxylation ($p=0.65$) and ribose-5-phosphate
($p=0.89$) do not meet $p<0.05$. Mouse HSPCs meet $p<0.05$ for HMP shunt ($p=0.0001$).
The other six mouse tasks do not, including glycolysis ($p=0.48$) and Complex~I
($p=0.34$).

Ni{\~n}o's G6PD/PPP fitness map is why HMP shunt is worth a cheap G6PD
assay~\cite{Nino2026TETmetabolism}. The permutation did not test G6PD activity.

Gene embedding clouds are the interpretability for that map (Fig.~\ref{fig:mice-genes},
Fig.~\ref{fig:human-genes}). Each point is one task gene. $x$ is the challenge main on
frozen gene embeddings. $y$ is the same interaction as the 2$\times$2 panels. Labels
collapse Complex~I subunits to Nduf / mt-Nd. Most effects coincide, with Complex~I as a
pile at the origin. Sdha and Ogdh map to succinate and $\alpha$-KG, Pkm maps to lactate,
and Tkt maps to NADPH and pentose. Those embeddings name genes off the pile. They are
not the permutation test. Challenge is not the same (IL-1 versus LPS), so the maps are not
stacked. These libraries contain HSPCs only. Human TET2 myeloid daughters amplify
inflammatory signalling while HSPCs blunt AP-1
adaptation~\cite{HuergaEncabo2026TET2inflammation}. The metabolomics extension stays in
HSPCs so the LC-MS matches the VNN.

\begin{figure}[p]
    \centering
    \includegraphics[width=\textwidth,height=0.82\textheight,keepaspectratio]{mice.png}
    \caption{\textbf{Mice.} McClatchy GSE209994 HSPC visible neural network (VNN),
        trained for 100 epochs on GPU. Each panel is the frozen task-encoding
        2$\times$2. Colour is relative mean. The $p$ under the panel is a competitive
        mean permutation of the gene-embedding interaction (9{,}999 gene shuffles;
        smallest possible $p=0.0001$). Bold is HMP shunt ($p=0.0001$). Panel titles
        match Fig.~\ref{fig:mice-genes}.}
    \label{fig:mice}
\end{figure}

\begin{figure}[p]
    \centering
    \includegraphics[width=\textwidth,height=0.82\textheight,keepaspectratio]{mice_genes.png}
    \caption{\textbf{Mice, gene embeddings.} Each point is a gene in that metabolic
        task. $x$ is the IL-1 main on frozen gene embeddings. $y$ is the same interaction as
        Fig.~\ref{fig:mice}. Both axes are normalized to $[-1,1]$. Labels are farthest from
        the origin (Nduf / mt-Nd collapsed).}
    \label{fig:mice-genes}
\end{figure}

\begin{figure}[p]
    \centering
    \includegraphics[width=\textwidth,height=0.82\textheight,keepaspectratio]{human.png}
    \caption{\textbf{Human.} GSE285379 FACS-HSPC VNN, TET2 $\times$ LPS (CTRL/LPS, 4
        libraries), trained for 100 epochs on GPU. Panels are encoding 2$\times$2, as in
        Fig.~\ref{fig:mice}. The $p$ under each panel is a competitive mean permutation
        of the gene-embedding interaction (9{,}999 gene shuffles; smallest $p=0.0001$).
        Bold includes glycolysis, Complex~I and Complex~II (all $p=0.0001$), HMP shunt
        ($p=0.0014$), Krebs-cycle NADH generation ($p=0.006$), and ATP regeneration
        ($p=0.043$). Pyruvate decarboxylation does not ($p=0.65$). Of the polar acids,
        succinate and $\alpha$-KG sit in TCA/OXPHOS tasks; lactate sits in glycolysis.
        Four libraries cannot decide an encoding interaction. Not merged into the mouse
        map.}
    \label{fig:human}
\end{figure}

\begin{figure}[p]
    \centering
    \includegraphics[width=\textwidth,height=0.82\textheight,keepaspectratio]{human_genes.png}
    \caption{\textbf{Human, gene embeddings.} Same axes as Fig.~\ref{fig:mice-genes}
        with LPS as the challenge main, normalized to $[-1,1]$. Four libraries cannot decide
        an interaction.}
    \label{fig:human-genes}
\end{figure}
\FloatBarrier

\section*{Proposed experiment}

If LC-MS shows succinate in the medium, a SUCNR1 reporter, or an osteoclast/HSPC
co-culture fed that medium, tests whether the cargo is
sensed~\cite{Guo2017Succinate,Cuminetti2026SUCNR1}. Osteoclast SUCNR1 is the
bone-remodeling path and HSPC SUCNR1 is the clone-size path. Medium that contains
succinate but does not move either assay leaves the niche sentence untested. If LC-MS
shows an $\alpha$-KG/succinate shift and 5hmC did not move at the cheap step, the
metabolites are not reaching catalysis. Kim already measured the myeloid lipidome on the
OGT-ACLY circuit. This polar panel, after a cheaper de-risking screen, is the
HSPC $\times$ IL-1 measurement Kim did not run.

Matched RNA-seq of the same cultures would later feed a multi-omics factor model and a
gene-regulatory network from this HSPC 2$\times$2. Those analyses are not required to
start the cheap assays or the LC-MS. Fracture-union prediction and IL-1$\beta$ blockade
remain cohort questions. This 2$\times$2 measures what polar cargo Tet2 HSPCs put into
the niche those cohorts already inhabit.

\section*{Conclusion}

Mouse and human HSPC RNA already indicate between-group differences at a short
gene-metabolite list. Most metabolic-task genes share one effect, so the encoding
2$\times$2 is not a group test. Mouse and human each get their own VNN, trained for
100 epochs on GPU. The $p$ under each panel is a competitive mean permutation on frozen
gene embeddings, with 9{,}999 gene shuffles (smallest possible $p=0.0001$). Mouse HSPCs
meet $p<0.05$ for HMP shunt. Human HSPCs meet it for glycolysis, both OXPHOS complexes,
HMP shunt, Krebs-cycle NADH generation, and ATP regeneration. Kim named the TET2 myeloid
lipid circuit. The literature has not measured the polar metabolome on this Tet2
$\times$ IL-1 HSPC 2$\times$2. TCA intermediates in HSPCs are metabolically quieter than
in downstream progenitors. Cheap CFU, 5hmC, and G6PD or Seahorse assays on the same four
arms come first, to de-risk the spend. Those assays test whether the cargo is used for
output, catalysis, and PPP fitness. Polar LC-MS of cells and medium is still planned. A
later SUCNR1 reporter or osteoclast/HSPC co-culture would test whether medium succinate
is sensed.

That niche is marrow inflammaging, with chronic myeloid demand and the cytokine tone
that expands the clone. Frailty is lost reserve in the osteoimmune niche that shares
those signals. Polar cargo from Tet2 HSPCs sits next to Kim's myeloid lipidome.
TET2-selective ligands for osteoclast SUCNR1 and HSPC SUCNR1 test for inflammasome tone
and resolution. Panel $p$ values come from the gene-set mean shuffles. Succinate and
$\alpha$-KG sit in TCA/OXPHOS tasks that meet $p<0.05$ in human, not in mouse. Lactate
sits in glycolysis, which meets that cutoff in human only. Mouse HMP shunt meets
$p<0.05$; mouse glycolysis and both OXPHOS complexes do not. Cheap cell assays de-risk
whether those scores are used. Polar LC-MS is still planned, to identify the metabolites
those scores point to.

\bibliographystyle{unsrtnat}
\bibliography{citations}

\end{document}
